FER+ established distribution-based supervision for multiple annotations of FER2013 images \cite{barsoum2016ferplus}. Washington et al. compared hard and soft targets on CAFE and showed that categorical performance and agreement with response distributions can favor different models \cite{washington2021crowdsourcing}. Other FER methods use disagreement through weighting and relabeling, uncertainty-aware target construction, or soft-label video augmentation \cite{wang2020scn,le2023uncertainty,kawamura2024midas}. We retain the observed votes and compare their training representations under matched conditions, rather than benchmark the additional mechanisms of those methods.

Selective classification trades accepted coverage against prediction risk \cite{geifman2017selective}; SelectiveNet jointly learns prediction and rejection \cite{geifman2019selectivenet}. Crowd-Calibrator estimates annotator distributions to inform rejection in subjective language tasks \cite{khurana2024crowd}. In its hate-speech comparison, soft-label models improve soft and calibration measures without corresponding categorical gains, whereas soft-label MaxProb performs worse on AUC/AUROC selection measures. Thus, disagreement between distributional and selective outcomes is not a new general finding. Our FER evaluation asks what remains after matching dominant target probabilities and ties, holding accepted images fixed, and accounting for the complete annotation vocabulary. We do not implement its learned crowd estimator or claim directly comparable benchmark scores.

Xia et al. show that conventional label smoothing can worsen selective classification by changing the ranking of correct and incorrect predictions, and examine post hoc logit normalization \cite{xia2025smoothing}. Our Uniform control instead uses image-dependent dominant probabilities and our loss scores disagreement with empirical votes, rather than binary correctness. Nevertheless, the ranking explanation is directly relevant. We therefore examine probability-based scores and a normalization sensitivity without retraining the existing models; these analyses assess ranking dependence rather than establish a new selector.

Calibration against a majority label also differs from agreement with annotator distributions. Baan et al. demonstrate this distinction and develop measures of human-judgment distribution properties \cite{baan2022calibration}; Guo et al. examine conventional confidence calibration \cite{guo2017calibration}. We report distributional distance, selected categorical disagreement and achieved calibration-transfer risk separately. Lower squared distance is not interpreted as a deployment risk guarantee.

The backbones and channel gates use established Swin, ResNet and squeeze-and-excitation designs \cite{liu2021swin,he2016resnet,hu2018senet}. Vats and Chadha describe Swin with squeeze-and-excitation for FER (unpublished preprint, 2023; \href{https://arxiv.org/abs/2301.10906}{arXiv:2301.10906}). Recent unpublished preprints examine ambiguity-aware calibration (Tao et al., 2026; \href{https://arxiv.org/abs/2603.22879}{arXiv:2603.22879}), dynamic disagreement prediction (Wang et al., 2026; \href{https://arxiv.org/abs/2609.17130}{arXiv:2609.17130}), and routing between ambiguity and distribution shift (Inoshita and Ueno, 2026; \href{https://arxiv.org/abs/2606.22725}{arXiv:2606.22725}). These related methods are not implemented baselines; their results do not validate ours. Our contribution concerns controlled target comparisons and their interpretation, with no architecture or rejection-method novelty claimed.
